\section{General symmetric kernels, complementary cuts, and odd orders}
\label{sec:generalkernels}
The bipartite hypothesis can be removed when the negative spectrum
is also separated from \(-1\).
The essential new transfer is a dependent, complementary-cut
estimate. Independent row and column selectors cannot simply be
conditioned to be complements: that event can have exponentially
small probability. The argument below keeps their dependence exactly.

\begin{lemma}[Complementary restriction with the full covariance retained]
\label{lem:complementaryrestriction}
Let \(N\geq2\) and let \(E=E^\top\) be an \(N\)-by-\(N\) matrix with
\(\|E\|\leq1\) and all squared column norms at most \(\alpha\geq0\).
Let \(R\) select each coordinate independently with probability
\(1/2\), and put \(R'=I-R\).
For every fixed \(A>0\), with probability \(1-N^{-A}\),
\[
 \|RER'\|^2\leq\tfrac14\|E\|^2+
       C_A\bigl(\sqrt{\alpha\log(2N)}+\alpha\log(2N)\bigr).
\]
The assertion remains true with polynomially high probability
after conditioning \(|R|=N/2\), when \(N\) is even, or after
also prescribing one coordinate in each part.
\end{lemma}
\begin{proof}
For columns \(c_j\) of \(E\), write
\[
 F=ERE-\tfrac12 E^2
       =\sum_j(\one_{j\in R}-1/2)c_jc_j^\top .
\]
The summands have norm at most \(\alpha\), and their matrix
variance is at most
\[
 \tfrac14\sum_j\|c_j\|_2^2c_jc_j^\top
              \preceq\tfrac{\alpha}{4}E^2 .
\]
Here is a finite signed-moment proof of the required tail.
Put \(A_j=c_jc_j^\top/2\), write
\(F_\epsilon=\sum_{j=1}^N\epsilon_jA_j\), and let
\(M_{2k}=2^{-N}\sum_{\epsilon\in\{-1,1\}^N}
                   \operatorname{tr}F_\epsilon^{2k}\).
For \(\epsilon^{(j)}\) obtained by flipping coordinate \(j\),
the exact identities are
\[
 \Delta_j=F_\epsilon-F_{\epsilon^{(j)}}=2\epsilon_jA_j,
 \qquad \Delta_j^2=4A_j^2,\qquad \sum_{j=1}^N\Delta_j=2F_\epsilon .
\]
The flip is an involution of the whole sign cube. Hence
\[
 4M_{2k}=2^{-N}\sum_{\epsilon\in\{-1,1\}^N}\sum_{j=1}^N
 \operatorname{tr}\!\left(
 \Delta_j(F_\epsilon^{2k-1}-F_{\epsilon^{(j)}}^{2k-1})\right).
\]
The trace mean-value inequality proved below retains both
powers and every mixed term. Apply it before the flip
bijection combines the two power contributions, giving
\[
 M_{2k}\leq(2k-1)\,2^{-N}\sum_{\epsilon\in\{-1,1\}^N}
                 \operatorname{tr}(V F_\epsilon^{2k-2}),
 \qquad V=\sum_{j=1}^NA_j^2 .
\]
In full, \(V=\tfrac14E\operatorname{diag}
(\|c_j\|_2^2)E^\top\). For every symmetric \(Z\) and
even \(q\), the diagonals of \(E^\top Z^qE\) are
nonnegative: for \(q=2r\) this is the Gram matrix
\((Z^rE)^\top(Z^rE)\).
Cyclic trace therefore gives
\[
 \operatorname{tr}(VZ^q)
 \leq\tfrac{\alpha}{4}\operatorname{tr}(E^2Z^q)
 \leq\tfrac{\alpha}{4}\operatorname{tr}(Z^q).
\]
For the last inequality, use the independent eigenbasis
overlap squares of \(E,Z\), the bounds
\(\lambda_i(E)^2\leq1\) and \(\lambda_j(Z)^q\geq0\).
Consequently
\[
 M_{2k}\leq N(2k-1)!!(\alpha/4)^k
          \leq N(\alpha k/2)^k .
\]
The even trace controls \(\|F_\epsilon\|^{2k}\), so, for
every \(k\geq1,t>0\),
\[
 2^{-N}\#\{\epsilon:\|F_\epsilon\|>t\}
          \leq\frac{N(\alpha k/2)^k}{t^{2k}}.
\]
For \(N\geq2,A>0,\alpha>0\), take
\(k=\lceil(A+4)\log(2N)\rceil\).
At \(t=\sqrt{2\alpha k}\) the bound is \(N4^{-k}\).
The inequalities \(\log4\geq1\) and
\(k\leq(A+5)\log(2N)\) give the explicit uniform form
\[
 \Pr\!\left\{\|F\|>
       \sqrt{2(A+5)\alpha\log(2N)}\right\}
             \leq N^{-A-3}.
\]
If \(\alpha=0\), every column vanishes and \(E=F=0\),
so the same conclusion holds. This proves the asserted
tail, with room for the conditioning costs below.
Crucially the other restriction has the exact relation
\[
 ER'E=\tfrac12 E^2-F .
\]
Consequently, retaining the same covariance \(F\),
\[
 \begin{split}
 \|RER'\|^2
 &=\|R'ERER'\|\\
 &\leq\tfrac12\|ER'\|^2+\|F\|\\
 &\leq\tfrac14\|E\|^2+\tfrac32\|F\|.
 \end{split}
\]
This proves the assertion without independent complementary
selectors or a discarded cross term.

The balanced event has probability at least \(1/(N+1)\).
With prescribed distinct coordinates \(a\in R,b\notin R\),
its probability is
\(2^{-N}\binom{N-2}{N/2-1}\geq1/(4(N-1))\).
Choosing the initial failure exponent larger absorbs these
polynomial conditioning costs.
\end{proof}

\paragraph{Maintained full complementary covariance and actual cut bound.}
The modules \texttt{HamiltonianComplementaryCovariance} and
\texttt{HamiltonianComplementaryCut} now prove the deterministic
part on the actual coordinate supports.
For any coordinate set \(D\), let \(R\) and \(S\) be its two
complementary diagonal projections, and put
\(F=ERE-\tfrac12E^2\).
The exact identities and bound are
\[
 ESE=\tfrac12E^2-F,\qquad F=F^\top,\qquad
 \|E[D,D^c]\|^2\leq\tfrac14\|E\|^2+\tfrac32\|F\|.
\]
The same \(F\) occurs with both signs before the operator estimates;
no independent-selector assumption is used.
The formal proof includes projection symmetry, idempotence,
complementarity and norm, the exact retained matrix, and its
dimension-preserving injective restriction.
Every column remains in the full signed covariance entry:
\[
 F_{ij}=\sum_{k=1}^N
       \bigl(\mathbf1_{k\in D}-\tfrac12\bigr)E_{ik}E_{kj}.
\]
It also proves the actual spectral consumer:
if \(0<\delta\leq1\), \(\|E\|\leq1-\delta\) and
\(\|F\|\leq\delta/128\), then
\[
 2\|E[D,D^c]\|\leq1-\tfrac34\delta .
\]
All eleven new deterministic covariance, coordinate-support and
gap declarations pass their combined focused foundational-axiom audit.
These statements are uniform over the finite matrix, its dimension,
and every cut; they do not restrict \(N\alpha\) to a sublogarithmic
range. The uniform probability tail for \(F\) is now formalized below.
Balanced and prescribed-coordinate conditioning in
Lemma~\ref{lem:complementaryrestriction} are now formalized below.
The full wide-cap entropy theorem is not yet claimed as formalized.

\paragraph{Maintained signed trace moment input.}
The three modules \texttt{HamiltonianTraceMeanValueScalar},
\texttt{HamiltonianMixedTraceReadout}, and
\texttt{HamiltonianTraceMeanValue} now prove the complete
trace mean-value inequality for all real symmetric matrices
\(X,Y\) and every integer \(k\geq1\):
\[
 \operatorname{tr}\!\left((X-Y)(X^{2k-1}-Y^{2k-1})\right)
 \leq \frac{2k-1}{2}
       \operatorname{tr}\!\left((X-Y)^2(X^{2k-2}+Y^{2k-2})\right).
\]
The scalar proof pairs every exponent of the exact difference
quotient with its reflection. Independent orthogonal
diagonalizations of \(X,Y\) then express both full trace
polynomials through the same nonnegative overlap squares
\((U^\top V)_{ij}^2\). Zero powers are included, so the pure
terms and every mixed compensation term retain exactly the
same finite support. No assumption on the signs of the
eigenvalues or on their multiplicities is imposed.
All seventeen new scalar, mixed-readout and trace declarations
compile and pass their focused foundational-axiom audit.

For the covariance above, put \(A_j=c_jc_j^\top/2\) and
\(F=\sum_{j=1}^N\epsilon_jA_j\), with independent uniform signs.
The compiled finite sign-flip identity yields the exact
variance calculation
\[
 V=\sum_j A_j^2
  =\frac14\sum_j\|c_j\|_2^2c_jc_j^\top
  \preceq \frac{\alpha}{4}E^2,
 \qquad
 \mathbb E\operatorname{tr} F^{2k}
 \leq (2k-1)\frac{\alpha}{4}\mathbb E\operatorname{tr} F^{2k-2}.
\]
These finite sign-flip moment and variance statements are now
Lean theorems, together with their uniform probability tail.
The eight maintained owners from \texttt{HamiltonianSignCube}
through \texttt{HamiltonianComplementaryPolynomialTail}
include the exact involutions, squared increments, full
variance factorization, all-order moments, actual norm
readout, the zero-variance case, and an exact bijection
between sign vectors and every coordinate cut.
All forty-eight new declarations compile and pass their
combined focused foundational-axiom audit.

In particular, their actual covariance tail is
\[
 \frac{\#\{D\subseteq[N]:
   \|ERE-\tfrac12E^2\|>
       \sqrt{2(A+5)\alpha\log(2N)}\}}{2^N}
          \leq N^{-A-3}.
\]
The same finite count bound holds for cuts violating the
actual quarter estimate
\[
 \|E[D,D^c]\|^2
 \leq\tfrac14\|E\|^2+
          \tfrac32\sqrt{2(A+5)\alpha\log(2N)}.
\]
The exact fixed-gap failure inclusion is also compiled.
These are uniform over every finite dimension \(N\geq2\)
and every \(\alpha\geq0\), with no small-\(N\alpha\)
restriction. Balanced and prescribed-coordinate conditioning
are also formalized below. The full wide-cap entropy
headline remains unfinished.

\paragraph{Maintained balanced and prescribed finite laws.}
The modules \texttt{HamiltonianBalancedCutCounts},
\texttt{HamiltonianPrescribedCutCounts},
\texttt{HamiltonianFiniteConditioning}, and
\texttt{HamiltonianConditionedComplementaryTail} complete
the conditioning step.
For \(N=2m\), \(m\geq1\), write
\[
 \mathcal B_m=\{\epsilon\in\{-1,1\}^N:|D_\epsilon|=m\},
 \quad
 \mathcal B_m(a,b)=\{\epsilon\in\mathcal B_m:
                           a\in D_\epsilon,\ b\notin D_\epsilon\}.
\]
For every distinct \(a,b\), their exact counts and budgets are
\[
 \begin{aligned}
 |\mathcal B_m|&=\binom Nm,&
 2^N&\leq(N+1)|\mathcal B_m|,\\
 |\mathcal B_m(a,b)|&=\binom{N-2}{m-1},&
 2^N&\leq4(N-1)|\mathcal B_m(a,b)|.
 \end{aligned}
\]
The two prescribed-coordinate count is a genuine
insertion/deletion bijection on the finite supports.
Both budgets follow from
\(2^{2r}\leq(2r+1)\binom{2r}{r}\), proved by summing
every binomial coefficient and comparing it to the middle one.
No endpoint case or coinciding-coordinate exception is suppressed;
the distinctness hypothesis is explicit, and \(N=2\) is included.

For any exceptional subset \(\mathcal E\) of the full sign cube
with \(|\mathcal E|/2^N\leq N^{-A-3}\), the exact
intersection counts give
\[
 \frac{|\mathcal E\cap\mathcal B_m|}{|\mathcal B_m|}
       \leq N^{-A},
 \qquad
 \frac{|\mathcal E\cap\mathcal B_m(a,b)|}{|\mathcal B_m(a,b)|}
       \leq N^{-A}.
\]
The denominators are proved positive from their count budgets.
The final factor absorption uses
\(N+1\leq N^3\) and \(4(N-1)\leq N^3\) for \(N\geq2\).
Applying these compiled statements to both the actual
covariance exceptional set and the actual coordinate
quarter-bound exceptional set proves the full finite
uniform-cut form of Lemma~\ref{lem:complementaryrestriction},
including both conditioning clauses.
Its stronger explicit error is
\[
 \tfrac32\sqrt{2(A+5)\alpha\log(2N)}.
\]
All eighteen new finite-count and conditional-tail declarations
compile and pass their focused foundational-axiom audit.
Weighted row concentration, the repaired balanced kernels and
the full general weighted entropy count remain separate
unfinished formal obligations.

\begin{lemma}[Scaling repair without a parity kernel]
\label{lem:aperiodicrepair}
Let \(Q\) be symmetric nonnegative of any order, and put
\[
 d=Q\one,\quad D=\diag(d),\quad
 \|d-\one\|_\infty\leq\varepsilon<1/2,\quad
 S=D^{-1/2}QD^{-1/2}.
\]
Assume \(S\) has simple eigenvalue \(1\) and all other
eigenvalues of modulus at most \(1-\delta\).
Then the scaling conclusion (5) holds. No bipartition is needed.
For sufficiently small \(\varepsilon L^2\), the repaired matrix
has nontrivial gap at least \(\delta/2\) and largest entry at most
twice that of \(Q\).
\end{lemma}
\begin{proof}
Use \(D,T,\pi\) as in Lemma~\ref{lem:scaling}, and let
\(\Pi_+f=\one\sum_i\pi_i f_i\), \(\Pi_0=I-\Pi_+\).
The same reversible spectral calculation gives
\[
 \|T^j\Pi_0\|_{\infty\to\infty}
       \leq C\min(1,\sqrt N(1-\delta)^j).
\]
There is no eigenvalue \(-1\), so \(D+Q=D(I+T)\) is invertible
on the whole space. Its inverse is specified by
\[
 (D+Q)^{-1}r
   =\tfrac12\Pi_+(D^{-1}r)
      +\sum_{j\geq0}(-1)^jT^j\Pi_0(D^{-1}r).
\]
Its infinity norm is \(O_\delta(L)\).
Equation (8) therefore has a contraction solution on the full
ball \(\|z\|_\infty\leq C_\delta\varepsilon L\);
the nonlinear term and its Lipschitz cost are exactly those
already calculated there. No orthogonality or gauge restriction
is imposed or required.
All entries of \(u=\one+z\) are positive.
The complete expansion \(ZQ+QZ+ZQZ\) gives the operator bound,
and the multiplier bound gives the entry bound.
Comparison with the normalized input and Weyl's inequality
give the stated gap after decreasing the admissible constant.
\end{proof}

\begin{lemma}[Most balanced cuts admit controlled stochastic scaling]
\label{lem:balancedcutscaling}
Let \(N\) be even and \(P=P^\top\) be nonnegative stochastic, with
\[
 \|P-J_N/N\|\leq1-\delta,\qquad \max_{ij}P_{ij}=\alpha,
 \qquad \alpha(\log(2N))^5\to0 .
\]
A proportion \(1-o(1)\) of all balanced ordered cuts \(X,Y\)
have the following property.
The symmetric bipartite kernel \(Q\) with cross block
\(2P[X,Y]\) has an exact stochastic scaling
\[
 P^{X,Y}=\diag(u)Q\diag(u),\qquad
 \|u-\one\|_\infty=O_\delta(\sqrt{\alpha}L^{3/2}).
\]
It has fixed nontrivial gap at least \(\delta/8\) and largest
entry at most \(4\alpha\).
The same conclusion holds among cuts with any fixed
\(a\in X,b\in Y\), uniformly in distinct \(a,b\).
\end{lemma}
\begin{proof}
Put \(E=P-J_N/N\).
Since \(P\) is symmetric stochastic,
\[
 \sum_j E_{ij}^2=\sum_jP_{ij}^2-1/N\leq\alpha .
\]
Lemma~\ref{lem:complementaryrestriction} thus applies.
For a balanced cut, exactly
\[
 2P[X,Y]-J_{N/2}/(N/2)=2E[X,Y].
\]
Its norm is at most
\(\sqrt{(1-\delta)^2+O(\sqrt{\alpha L}+\alpha L)}\),
hence at most \(1-3\delta/4\) for sufficiently large \(N\).
This bounds the second singular value of the cross block
by the rank-one approximation inequality.

Scalar concentration of each weighted row sum, with variance
at most \(\alpha/4\), gives
\[
 \|Q\one-\one\|_\infty
             \leq C(\sqrt{\alpha L}+\alpha L).
\]
Union bounds and the same polynomial conditioning costs give
this simultaneously with the operator estimate.
They also apply if \(P\) has nonzero diagonal entries:
every cross row sum is an actual selected sum, and each
diagonal contribution is at most \(\alpha\).
Normalizing the degrees changes the symmetric operator by
\(O(\sqrt{\alpha L}+\alpha L)\).
The normalized bipartite operator has its exact eigenvalues
\(1,-1\), and the remaining spectrum stays a fixed distance
from both; simplicity follows from the second-singular-value
bound.
Now \(\varepsilon L^2=O(\sqrt\alpha L^{5/2})=o(1)\).
Apply Lemma~\ref{lem:scaling}.
The multiplier, maximum-entry and gap conclusions follow,
with constants uniform over the good cuts and endpoints.
\end{proof}

\begin{theorem}[Hamiltonian-weight universality at every order]
\label{thm:generalweightedentropy}
Fix \(0<\delta<1,C>0\), and let \(K_{\rm w}=3\cdot10^7\).
Let \(P=P_N\) be any symmetric nonnegative stochastic kernels,
with
\[
 \|P-J_N/N\|\leq1-\delta,\qquad
 \alpha=\max_{ij}P_{ij}\leq C(\log(2N))^{-K_{\rm w}}.
\]
Loops are permitted in the kernel; Hamilton paths and cycles
use distinct vertices and no loops.
For all sufficiently large \(N\), of either parity,
\[
 \begin{split}
 \log Z_{\rm H}(P)
       &=-N+O_{\delta,C}(N\alpha^{1/8}\log(2N)),\\
 \log Z_{\rm path}(P;a,b)
       &=-N+O_{\delta,C}(N\alpha^{1/8}\log(2N))
                         \quad(a\ne b).
 \end{split}
\]
The endpoint error and threshold are uniform over all \(a,b\).
Together with the bipartite theorems, this covers both a fixed
aperiodic spectral gap and the exact balanced parity case.
It makes no assertion uniform in a spectral gap tending to zero.
\end{theorem}
\begin{proof}
The permanent upper bounds, including path minors, are already
Lemma~\ref{lem:diffusepermanent} and Lemma~\ref{lem:pathminor}.
Write \(R_N=N\alpha^{1/8}\log(2N)\).
Since \(\alpha\geq1/N\), this absorbs \(O(\log N)\).

Suppose first that \(N\) is even.
For every good balanced cut apply the bipartite weighted
theorem to \(P^{X,Y}\).
For a Hamilton cycle crossing that cut, the exact relation is
\[
 \prod_{ij\in E(C)}P^{X,Y}_{ij}
       =2^N\left(\prod_vu_v^2\right)
                         \prod_{ij\in E(C)}P_{ij}.
\]
The logarithm of the multiplier is
\(O_\delta(N\sqrt\alpha L^{3/2})=o(R_N)\).
The weighted theorem is uniform on these cuts, since their
gap is at least \(\delta/8\) and their largest entry at most
\(4\alpha\). Each good cut therefore contributes at least
\[
 2^{-N}\exp(-N-O_{\delta,C}(R_N))
\]
to the Hamilton sum in the original kernel.

Every even Hamilton cycle has exactly two balanced ordered
alternating cuts. Hence the exact double-counting identity is
\[
 2Z_{\rm H}(P)
       =\sum_{\substack{X\subseteq[1,N]\\|X|=N/2}}
                        Z_{\rm H}(P\hbox{ restricted to }X,X^c).
\]
There are \(\binom N{N/2}\) cuts and a proportion \(1-o(1)\)
are good. The factor \(2^{-N}\) cancels their exponential
number, leaving only \(O(\log N)\) in the logarithm.
This proves the cycle lower bound at even orders.

For a fixed oriented Hamilton path from \(a\) to \(b\),
there is exactly one balanced cut with \(a\in X,b\in Y\)
which it crosses throughout.
The exact scaling factor for its weight is instead
\[
 2^{N-1}\frac{\prod_vu_v^2}{u_a u_b}.
\]
Apply Theorem~\ref{thm:bipartitepathentropy} on each good cut
with these endpoints.
There are \(\binom{N-2}{N/2-1}\) such cuts, almost all good
uniformly in \(a,b\). Their exponential number cancels
\(2^{-(N-1)}\), again up to a polynomial factor.
This proves the even-order path lower bound, including all
endpoint and scaling factors.

Here is the reduction for odd orders.
Remove an arbitrary vertex \(v\), leaving the principal
matrix \(Q\) of order \(n=N-1\).
Its row defect is at most \(\alpha\).
Interlacing gives
\[
 \lambda_2(Q)\leq1-\delta,\qquad
 \lambda_{\min}(Q)\geq-1+\delta .
\]
Its normalized matrix differs from \(Q\) by \(O(\alpha)\),
has exact eigenvalue 1, and all other eigenvalues have modulus
at most \(1-\delta/2\), for large \(N\).
Lemma~\ref{lem:aperiodicrepair} therefore supplies
\[
 \widetilde P=\diag(u)Q\diag(u),\quad
 \|u-\one\|_\infty=O_\delta(\alpha L),\quad
 \max\widetilde P_{ij}\leq2\alpha
\]
with a fixed positive aperiodic gap.
The even-order path conclusion just proved applies to it.
Transferring a path back to \(Q\) uses the exact factor
\(\prod u_i^2/(u_a u_b)\), costing \(O_\delta(N\alpha L)=o(R_N)\).

Choose two distinct off-diagonal neighbours \(a,b\) of \(v\)
with \(P_{va},P_{vb}\geq1/(2N)\).
Such neighbours exist for large \(N\): after excluding the
diagonal and the largest off-diagonal entry, the remaining
row mass is at least \(1-2\alpha\).
Close a spanning \(a\)-to-\(b\) path of \(Q\) through \(v\).
This is an injective construction of Hamilton cycles, and
\[
 Z_{\rm H}(P)\geq
 P_{va}P_{vb}\,Z_{\rm path}(Q;a,b)
       \geq (4N^2)^{-1}\exp(-n-O_{\delta,C}(R_N)).
\]
Its additional logarithmic cost is \(O(\log N)\).
This proves the odd cycle bound without assuming a bipartition.

Finally, for prescribed endpoints \(a,b\) at odd order, remove
\(a\) and repair its principal matrix in the same way.
There is a neighbour \(c\notin\{a,b\}\) with
\(P_{ac}\geq1/(2N)\), since excluding the two specified entries
leaves mass at least \(1-2\alpha\).
Prepend \(a,c\) to a spanning path from \(c\) to \(b\) on the
remaining even-order matrix.
The construction is injective, its scaling cost is
\(O_\delta(N\alpha L)\), and its extra edge costs only
\(O(\log N)\).
This proves the final lower bound uniformly in the endpoints.
No parity case or diagonal mass has been omitted.
\end{proof}

\begin{corollary}[Sparse graph entropy without a bipartite restriction]
\label{cor:generalsinkhornentropy}
Let simple graphs \(G_N\), of any order, admit positive
symmetric stochastic scalings
\(P_N=\diag(u)A_{G_N}\diag(u)\).
Suppose either the aperiodic hypotheses of
Theorem~\ref{thm:generalweightedentropy}, or the balanced
bipartite hypotheses of Theorem~\ref{thm:weightedentropy}, hold.
Then
\[
 \log h(G_N)
   =2\max_{\substack{x_e\geq0\\\sum_{e\ni v}x_e=1}}
                \sum_{e\in E(G_N)}x_e\log(1/x_e)
      -N+O_{\delta,C}(N\alpha^{1/8}\log(2N)).
\]
Moreover the logarithmic cost of imposing the one-cycle
condition on weighted permutations is at most the same
order, for all orders covered here.
\end{corollary}
\begin{proof}
The cycle-weight multiplier is always \(\prod_vu_v^2\).
The entropy-maximization calculation in
Corollary~\ref{cor:sinkhornentropy} uses only the row equations,
so applies to general graphs as well.
The full permanent and Hamiltonian permanent change by the
same multiplier. Their comparison follows from the two
weighted upper and lower estimates just proved.
\end{proof}
